\documentclass[10pt,a4paper]{amsart}
\usepackage[margin=19mm]{geometry}
\usepackage{amsmath,amssymb,mathtools}
\usepackage[scr=rsfs,cal=euler]{mathalfa}
\usepackage{xfrac}
\usepackage[unicode,hidelinks,bookmarks=false]{hyperref}
\newcommand{\ie}{i.e.\ }
\newcommand{\cf}{cf.\ }
\newcommand{\resp}{respectively\ }
\newcommand{\Subsection}{\subsection}
\providecommand{\nobreakdash}{\nobreak\hbox{-}}
\setlength{\emergencystretch}{2em}
\multlinegap0pt
\usepackage{bm}
\usepackage{bbm}
\usepackage{dsfont}
\usepackage{xspace}
\usepackage{cancel}
\usepackage{mathtools}
\usepackage{amsthm}
\makeatletter
\@ifundefined{defn}{\newtheorem{defn}{Defn}}{}
\@ifundefined{claim}{\newtheorem{claim}[defn]{Claim}}{}
\makeatother
\title[Distinguishing colorings, proper colorings, and covering pro]{Distinguishing colorings, proper colorings, and covering properties without the Axiom of Choice}
\author{Amitayu Banerjee, Zal\'an Moln\'ar, Alexa Gopaulsingh}
\date{Source: arXiv:2309.06116v3 (CC BY 4.0)}
\begin{document}
\maketitle

\noindent\emph{Source and licence.} Distinguishing colorings, proper colorings, and covering properties without the Axiom of Choice. Version arXiv:2309.06116v3 (CC BY 4.0), distributed under the Creative Commons Attribution 4.0 International licence, as recorded in the arXiv licence metadata for this exact version and shown on the abstract page https://arxiv.org/abs/2309.06116v3. Full rights notice: licenses/arxiv-2309.06116-CC-BY-4.0.txt.\par\smallskip

\noindent\textit{Editorial scope.} Counted unit: the Claim whose source label is \texttt{None} in the article (source0.tex lines 538--540, source label \texttt{None}), delivered together with the article's own complete original proof of it (source0.tex lines 541--562, one \texttt{proof} environment of 22 lines). No mathematics is written, repaired or re-derived: statement and proof are the article's own text line for line. Required context retained uncounted: none. Every definition, display and label of those lines is retained unchanged. Omitted from this package and not redistributed: 1--537, 563--1448. Nothing inside a retained excerpt is omitted or altered: every retained body is delivered exactly as the article's lines are written, so no omission receipt of any kind accompanies this package. Citations: the retained excerpts carry no citation command at all, so no bibliography entry of the article is transcribed and no reference is redistributed. Reference adaptation: none was needed and none was made; every reference inside the delivered unit resolves inside it, so no unresolved reference or citation is produced. Printed slips kept literal: no printed word, symbol, hypothesis or number of the counted statement or of its proof was altered. Layout and class: the article's own class is \texttt{amsart} with packages including hyperref, enumerate, geometry, amsmath, amssymb, amsthm, amsfonts, graphicx, tikz; the wrapper is the lane's pinned \texttt{amsart} editorial wrapper at 10pt on A4, which restates the theorem-like environments the retained text needs and adds the article's own macro definitions that the retained text needs (none), with extra wrapper packages (bm, bbm, dsfont, xspace, cancel, mathtools). The delivered numbering is the unit's own. Assets: no retained span contains a figure environment, a graphics-inclusion command, a PDF-page inclusion, a table or a TikZ picture; no image file is copied into this package, the package declares no asset dependency and no image file is redistributed with it. Licence: version arXiv:2309.06116v3 of this article is distributed under the Creative Commons Attribution 4.0 International licence, as recorded in the arXiv licence metadata for this exact version and shown on the abstract page https://arxiv.org/abs/2309.06116v3. A full rights notice is delivered at licenses/arxiv-2309.06116-CC-BY-4.0.txt. Characters: every retained excerpt is pure ASCII, so the delivered engine, XeLaTeX with its default Latin Modern text font, reports no missing character.
		\begin{claim}
			{\em $t', t''$, and $t_m$ are fixed by every automorphism for each non-negative integer $m$.}
		\end{claim}

		\begin{proof}
			Fix non-negative integers $n,m,r$. The first-order $\mathcal{L}$-formula
			\[ \mathsf{Deg}_n(x) := \exists x_0\dots \exists x_{n-1} \big(\bigwedge_{i\neq j}^{n-1}x_i\neq x_j\wedge  \bigwedge_{i<n}x\neq x_i \wedge \bigwedge_{i<n} Exx_i \wedge \forall y(Exy \to \bigvee_{i<n}y=x_i)\big)\]
			expresses the property that a vertex $x$ has degree $n$, where $Eab$ denotes the existence of an edge between vertices $a$ and $b$.  We define the following first-order $\mathcal{L}$-formula:
			
			\begin{flushleft}
				$\varphi(x) := \mathsf{Deg}_1(x)\wedge \exists y(Exy \wedge \mathsf{Deg}_2(y))$.
				%$\psi_r(x,y) : =  \exists x_0\dots \exists x_r\big(x_0 = y \wedge x_r = x\wedge \bigwedge_{i\neq j}^r x_i\neq x_j \wedge \bigwedge_{i<r} Ex_ix_{i+1} \wedge\exists z(Exz\wedge z\neq x_{r-1})\big)$,
				%and $\varphi_n(x):=\exists y ( \mathsf{Deg}_2(y) \wedge \psi_n(x,y))$.\footnote{i.e., $\psi_r(x,y)$ depicts there exists a path of length $r$ from $y$ to $x$ such that $x$ has degree more than 1 and $\phi_{n}(x)$ depicts there exists a vertex of degree $2$ for which $\psi_{n}(x,y)$ holds.}
			\end{flushleft}
			
			It is easy to see the following:
			
			\begin{enumerate}
				\item[(i)] $t''$ is the unique vertex such that $H_1^1\models \varphi(t'')$. This means $t''$ is the unique vertex such that $\deg(t'')=1$ and $t''$ has a neighbor of degree $2$.
				\item[(ii)] $t'$ is the unique vertex such that $H_1^1\models  \mathsf{Deg}_2(t')$. So $t'$ is the unique vertex with $\deg(t')= 2$.
				%\item [(iii)] $t_m$ is the unique vertex such that $H_1^1\models \varphi_{m+1}(t_m)$. This means  $t_m$ is the unique vertex having path length $m-1$ from $t'$ and $\deg(t_{m})>1$. 
			\end{enumerate}
			Fix any automorphism $\tau$. Since every automorphism preserves the properties mentioned in (i)--(ii), $t'$ and $t''$ are fixed by $\tau$. 
			The vertices $t_{m}$ are fixed by $\tau$ by induction as follows: Since $t_{i}$ is the unique vertex of path length $i+1$ from $t''$ such that the degree of $t_{i}$ is greater than 1, where $i\in \{0,1\}$, we have that $t_{0}$ and $t_{1}$ are fixed by $\tau$.
			Assume that $\tau(t_{l})=t_{l}$ for all $l<m-1$. We show that $\tau(t_{m})=t_{m}$. Now, $\tau(t_{m})$ is a neighbour of $\tau(t_{m-1}) = t_{m-1}$ which is of degree at least $2$, so $\tau(t_{m})$ must be either $t_{m-2}$ or $t_{m}$, but $t_{m-2} = \tau(t_{m-2})$ is already taken. So, $\tau(t_{m})=t_{m}$.
		\end{proof}

\end{document}
